\chapter*{Notação}
\label{cap:notacao-front}
\addcontentsline{toc}{chapter}{Notação}
\markboth{Notação}{Notação}

A notação a seguir é usada de forma consistente em todo o livro. Onde um trabalho
original empregava símbolo diferente, ele foi harmonizado; a correspondência está
registrada no \cref{cap:notacao}.

\section*{Convenções gerais}

\begin{center}
\begin{tabular}{ll}
\toprule
\textbf{Forma} & \textbf{Significado} \\
\midrule
$x$, $\alpha$              & escalar \\
$\vect{x}$                 & vetor (coluna, em negrito) \\
$\mat{X}$, $\mat{\Sigma}$  & matriz (em negrito maiúsculo) \\
$\vect{x}^{\top}$          & transposto \\
$\mat{\Sigma}^{-1}$        & inversa \\
$\|\vect{x}\|$             & norma euclidiana \\
$\hat{y}$                  & valor predito \\
$\bar{x}$                  & média amostral \\
$\R^{d}$                   & espaço euclidiano de dimensão $d$ \\
\bottomrule
\end{tabular}
\end{center}

\section*{Dados e séries temporais}

\begin{center}
\begin{tabular}{ll}
\toprule
\textbf{Símbolo} & \textbf{Significado} \\
\midrule
$T$              & número de passos de tempo de uma janela \\
$S$              & número de sensores (canais) \\
$\mat{X} \in \R^{T \times S}$ & janela multivariada de sensores \\
$\vect{x}_t \in \R^{S}$       & leitura de todos os sensores no instante $t$ \\
$y \in \{0,1,\dots,9\}$       & rótulo de classe do 3W \\
$\mathcal{K}$    & conjunto de classes conhecidas \\
$\mathcal{U}$    & conjunto de classes desconhecidas (novidades) \\
$\vect{b}$       & janela de linha de base, anterior ao evento \\
$\vect{e}$       & janela de evento \\
\bottomrule
\end{tabular}
\end{center}

\section*{Modelos}

\begin{center}
\begin{tabular}{ll}
\toprule
\textbf{Símbolo} & \textbf{Significado} \\
\midrule
$f_{\theta}$     & modelo com parâmetros $\theta$ \\
$g_{\phi}$       & codificador (\ing{encoder}) do autoencoder \\
$h_{\psi}$       & decodificador (\ing{decoder}) do autoencoder \\
$\vect{z} = g_{\phi}(\mat{X}) \in \R^{d}$ & representação latente \\
$d$              & dimensão do espaço latente \\
$\hat{\mat{X}} = h_{\psi}(\vect{z})$ & reconstrução \\
$c_k$            & classificador binário da classe $k$ \\
$s_k \in [0,1]$  & escore do classificador binário $c_k$ \\
$\tau$           & limiar de decisão \\
\bottomrule
\end{tabular}
\end{center}

\section*{Estatística e avaliação}

\begin{center}
\begin{tabular}{ll}
\toprule
\textbf{Símbolo} & \textbf{Significado} \\
\midrule
$\mu$, $\sigma$    & média e desvio-padrão \\
$\mat{\Sigma}$     & matriz de covariância \\
$d_M(\vect{x})$    & distância de Mahalanobis \\
$\lambda$          & parâmetro de regularização da covariância \\
$k$                & multiplicador de desvio-padrão em limiares $\mu + k\sigma$ \\
$K$                & número de agrupamentos (\ing{clusters}) \\
VP, VN             & verdadeiros positivos e verdadeiros negativos \\
FP, FN             & falsos positivos e falsos negativos \\
$\text{ACC}$, $\text{ACC}_b$ & acurácia e acurácia balanceada \\
$\text{REC}$, $\text{PR}$, $\text{SP}$ & revocação, precisão e especificidade \\
$F_1$              & medida F1 \\
\bottomrule
\end{tabular}
\end{center}

\begin{destaque}[Colisão de símbolos]
Dois usos de $k$ convivem no livro e não podem ser confundidos: o \emph{multiplicador
de desvio-padrão} em limiares da forma $\mu + k\sigma$ (\cref{cap:dual},
\cref{cap:mundoaberto}) e o \emph{número de vizinhos ou de posições em um
\ing{ranking}} nas métricas ``top-$k$'' (\cref{cap:resultados-agente}). O contexto
sempre desambigua; quando houver risco, o significado é repetido por extenso.
Da mesma forma, $K$ maiúsculo designa exclusivamente o número de agrupamentos.
\end{destaque}
